\begin{frame}[standout]

    Avez-vous été attentifs ?

    À vous de jouer !

\end{frame}

{
\setbeamercovered{transparent}

\begin{frame}{Q1}

    \textbf{Quels sont les quatre principaux types de bases de données NoSQL ?}

    \begin{itemize}
        \item \uncover<1>{\textcolor{kh-red}{Colonne, Clé-Valeur, XML, Graph}}
        \item \uncover<1>{\textcolor{kh-blue}{JSON, Colonne, Document, Graph}}
        \item \uncover<1>{\textcolor{kh-green}{Clé-Valeur, Table, Document, SQL}}
        \item \textbf<2>{\textcolor{kh-yellow}{Clé-Valeur, Colonne, Document, Graph}}
    \end{itemize}

\end{frame}

\begin{frame}{Q2}

    \textbf{Dans une base de données de type document, sous quelle forme les données sont-elles stockées ?}

    \begin{itemize}
        \item \uncover<1>{\textcolor{kh-red}{Clé-valeur}}
        \item \textbf<2>{\textcolor{kh-blue}{JSON ou BSON}}
        \item \uncover<1>{\textcolor{kh-green}{Table}}
        \item \uncover<1>{\textcolor{kh-yellow}{Colonne}}
    \end{itemize}

\end{frame}

\begin{frame}{Q3}

    \textbf{Quelle base de données NoSQL est conçue pour les requêtes de graphe complexe ?}

    \begin{itemize}
        \item \uncover<1>{\textcolor{kh-red}{Redis}}
        \item \uncover<1>{\textcolor{kh-blue}{Apache Cassandra}}
        \item \textbf<2>{\textcolor{kh-green}{Neo4j}}
        \item \uncover<1>{\textcolor{kh-yellow}{MongoDB}}
    \end{itemize}

\end{frame}

\begin{frame}{Q4}

    \textbf{Quelle base de données est plus adaptée pour un jeu vidéo ?}

    En tant que membre d'une équipe de développeurs de jeux, vous avez été chargé de concevoir un moyen de suivre les possessions des joueurs (achetées dans un catalogue, échangées entre joueurs ou via des récompenses). Les possessions sont classées en vêtements, outils, livres et pièces de monnaie. Les joueurs peuvent avoir un nombre illimité de possessions de n'importe quel type. Les joueurs peuvent rechercher d'autres joueurs possédant des types de possessions particuliers pour des échanges. Le concepteur du jeu informe qu'il y aura probablement de nouveaux types de possessions et de nouvelles façons de les acquérir à l'avenir. Quel type de base de données recommanderiez-vous d'utiliser ?

    \begin{itemize}
        \item \uncover<1>{\textcolor{kh-red}{Base de données transactionnelle}}
        \item \uncover<1>{\textcolor{kh-blue}{Base de données orientée colonne}}
        \item \textbf<2>{\textcolor{kh-green}{Base de données orientée document}}
        \item \uncover<1>{\textcolor{kh-yellow}{Base de données analytique}}
    \end{itemize}

\end{frame}

\begin{frame}{Q5}

    \textbf{Quelle base de données est plus adaptée pour un logiciel ?}

    Un développeur de logiciels vous demande conseil sur le stockage des données. Il dispose de centaines de milliers d'objets JSON de 1 Ko qui doivent être accessibles en moins d'une milliseconde si possible. Tous les objets sont référencés par une clé. Il n'est pas nécessaire de rechercher des valeurs dans le contenu de la structure JSON. Quel type de base de données NoSQL recommanderiez-vous ?

    \begin{itemize}
        \item \uncover<1>{\textcolor{kh-red}{Base de données orientée graphe}}
        \item \uncover<1>{\textcolor{kh-blue}{Base de données analytique}}
        \item \uncover<1>{\textcolor{kh-green}{Base de données orientée colonne}}
        \item \textbf<2>{\textcolor{kh-yellow}{Base de données clé-valeur}}
    \end{itemize}

\end{frame}

\begin{frame}{Q6}

    \textbf{Quelle base de données est plus adaptée pour un logiciel ?}

    Une développeuse de logiciels vous demande conseil sur le stockage de données. Il a des centaines de milliers d'objets JSON de 10 Ko qui doivent pouvoir être recherchés par la plupart des attributs de la structure JSON. Quel type de base de données NoSQL recommanderiez-vous ?

    \begin{itemize}
        \item \textbf<2>{\textcolor{kh-red}{Base de données orientée document}}
        \item \uncover<1>{\textcolor{kh-blue}{Base de données orientée colonne}}
        \item \uncover<1>{\textcolor{kh-green}{Base de données analytique}}
        \item \uncover<1>{\textcolor{kh-yellow}{Base de données clé-valeur}}
    \end{itemize}

\end{frame}

%\begin{frame}{Q7}
%
%    \textbf{Une programmeuse conçoit une base de données pour prendre en charge les requêtes ad hoc. Quel type de schéma de données est-il susceptible d'utiliser ?}
%
%    \textit{Une requête ad hoc est une requête qui est créée pour répondre à un besoin d'information spécifique, souvent de manière ponctuelle, et qui n'est pas prévue ou récurrente dans les opérations normales de la base de données.}
%
%    \begin{itemize}
%        \item \uncover<1>{\textcolor{kh-red}{Online Transaction Processing (OLTP)}}
%        \item \textbf<2>{\textcolor{kh-blue}{Online Analytical Processing (OLAP)}}
%        \item \uncover<1>{\textcolor{kh-green}{Schéma normalisé}}
%        \item \uncover<1>{\textcolor{kh-yellow}{Schéma Dénormalisé}}
%        \item \uncover<1>{\textcolor{kh-teal}{Graphe de connaissances}}
%    \end{itemize}
%
%\end{frame}

%\begin{frame}{Q8}
%
%    \textbf{}
%
%    \begin{itemize}
%        \item \uncover<1>{\textcolor{kh-red}{}}
%        \item \uncover<1>{\textcolor{kh-blue}{}}
%        \item \textbf<2>{\textcolor{kh-green}{}}
%        \item \uncover<1>{\textcolor{kh-yellow}{}}
%    \end{itemize}
%
%\end{frame}
%
%\begin{frame}{Q9}
%
%    \textbf{}
%
%    \begin{itemize}
%        \item \uncover<1>{\textcolor{kh-red}{}}
%        \item \uncover<1>{\textcolor{kh-blue}{}}
%        \item \textbf<2>{\textcolor{kh-green}{}}
%        \item \uncover<1>{\textcolor{kh-yellow}{}}
%    \end{itemize}
%
%\end{frame}
%
%\begin{frame}{Q10}
%
%    \textbf{}
%
%    \begin{itemize}
%        \item \uncover<1>{\textcolor{kh-red}{}}
%        \item \uncover<1>{\textcolor{kh-blue}{}}
%        \item \textbf<2>{\textcolor{kh-green}{}}
%        \item \uncover<1>{\textcolor{kh-yellow}{}}
%    \end{itemize}
%
%\end{frame}

}
